\documentclass[12pt]{article}
\usepackage[utf8]{inputenc}
\usepackage[T1]{fontenc}
\usepackage[margin=1in]{geometry}
\usepackage{setspace}
\usepackage{xcolor}
\usepackage[colorlinks=true, urlcolor=blue!60!black, linkcolor=blue!60!black]{hyperref}
\setlength{\parindent}{0em}
\setlength{\parskip}{0.6em}
\renewcommand{\baselinestretch}{1.0}
\begin{document}
\begin{center}
{\LARGE \textit{Youngstown Sheet \& Tube Co. v. Sawyer}, 343 U.S. 579 (1952)\par}
\vspace{0.5em}
{\large Case Brief --- Two-Page Edition\par}
\end{center}

\section{Memory Jogger}
In the ``Steel Seizure Case,'' the Court held 6--3 that President Truman lacked statutory or constitutional power to seize strike-bound steel mills during the Korean War; Justice Jackson's three-zone framework became the test for executive power (\href{https://supreme.justia.com/cases/federal/us/343/579/}{343 U.S. 579}).

\section{Detailed Case Facts}
Facing a nationwide steel strike in April 1952, Truman declined to use the Taft-Hartley Act's cooling-off injunction or statutory condemnation and instead ordered Commerce Secretary Sawyer to seize and run the mills, citing only the Constitution, the laws, and his role as Commander in Chief. Congress took no action, and the companies sued.

\section{Procedural History}
Judge Pine enjoined the seizure; the D.C. Circuit stayed the injunction; the Supreme Court granted certiorari and affirmed within weeks.

\section{Judicial Votes}
Affirmed, 6--3. Black wrote for Frankfurter, Douglas, Jackson, and Burton, who each also concurred; Clark concurred in the judgment; Vinson, joined by Reed and Minton, dissented.

\section{Holding}
Presidential power to seize private property must come from Congress or the Constitution. No Article II clause authorizes seizing mills to avert a labor strike, a lawmaking function belonging to Congress, which had declined to grant seizure power.

\section{Analysis of Opinions}
Black reasoned that the Take Care Clause confines the President to recommending and vetoing laws and that Congress had rejected seizure power in 1947. Frankfurter allowed that long acquiesced-in practice can gloss executive power but found none here; Douglas stressed that only Congress can pay just compensation. Jackson's concurrence divided presidential power into three zones---``at its maximum'' with congressional authorization, a ``zone of twilight'' when Congress is silent, and ``at its lowest ebb'' when contrary to Congress's will---and placed the seizure in the third. Clark required the President to follow congressional procedures. Vinson, dissenting, invoked historical emergency actions and mocked a ``messenger-boy concept of the Office.''

\section{Examples: Future Applications}
\textbf{Same-Side:} seizing semiconductor plants to avert a strike, contrary to a labor statute, is at the ``lowest ebb.'' \textbf{Opposite-Side:} freezing foreign assets under IEEPA, as in \textit{Dames \& Moore}, is at maximum authority. \textbf{Fence-Sitter:} an emergency cybersecurity directive to private utilities amid congressional silence, in the ``zone of twilight.''

\section{Critique}
Formalists praise Black's defense of congressional lawmaking, while Goldsmith and Manning see a partial vindication of Vinson's ``completion power.'' Functionalists celebrate Jackson, though Mark Rosen questions Zone 3's assumption of ``categorical congressional supremacy.''

\section{Key Quotations}
\begin{quote}
``While the Constitution diffuses power the better to secure liberty, it also contemplates that practice will integrate the dispersed powers into a workable government. It enjoins upon its branches separateness but interdependence, autonomy but reciprocity. Presidential powers are not fixed but fluctuate depending upon their disjunction or conjunction with those of Congress.''
\end{quote}
\noindent --- Justice Jackson, concurring, 343 U.S. at 635.

\end{document}
